\section{Introduction}
\label{sec:introduction}

For a matrix $X=(x_{ij})$ of independent variables, the permanent is
\[
 \per_m(X)=\sum_{\sigma\in S_m}\prod_{i=1}^m x_{i,\sigma(i)}.
\]
The permanent--determinant problem asks how large a determinant of
affine-linear forms is needed to express $\per_m$. Valiant connected
determinant projections to arithmetic formulas: his construction represents
a polynomial computed by a formula with $s$ addition or multiplication
operations as a determinant of order $s+2$
\cite[Theorem~1]{Valiant1979Algebra}. His algebraic completeness theorem
shows that, in characteristic different from two, polynomial families
defined by summing polynomial-size formulas over Boolean strings of
polynomial length are projections of polynomial-size permanents
\cite[Theorem~2]{Valiant1979Algebra}. Valiant's projections substitute
individual variables or field constants. The exact model below allows
every matrix entry to be an arbitrary affine-linear form.

For $f\in\C[x_1,\ldots,x_N]$, its \emph{determinantal complexity}
$\dc(f)$ is the least positive integer $n$ for which
\[
 f=\det\left(A_0+\sum_{j=1}^N x_jA_j\right),
 \qquad A_0,\ldots,A_N\in\Mat_n(\C).
\]
Its \emph{border determinantal complexity} $\bdc(f)$ is the least positive
integer $n$ such that $f$ lies in the coefficientwise closure of
\begin{equation}
 \left\{\det\left(A_0+\sum_{j=1}^N x_jA_j\right):
 A_0,\ldots,A_N\in\Mat_n(\C)\right\}.
 \label{eq:border-definition}
\end{equation}
The closure is taken in the finite-dimensional space of polynomials of
degree at most $n$; a polynomial of larger degree cannot belong to it.
We allow $f$ to be constant or zero. Coefficientwise closure agrees with
Zariski closure here, as explained in Section~\ref{sec:models}. The size
$n$ stays fixed along an approximating sequence, while the complex
coefficients of its matrix entries are unrestricted.

\begin{theorem}
\label{thm:main}
There exist absolute constants $c>0$ and $m_0$ such that, for every integer
$m\ge m_0$,
\[
 \bdc(\per_m)\ge c m^3.
\]
One may take $c=1/(5\,529\,600e)$ and $m_0=1408$.
\end{theorem}

The constants are not optimized. Proposition~\ref{prop:model-equivalence}
identifies this coefficientwise model with the projective determinant-orbit
closure formulation, including equal permanent and determinant orders.

The same numerical lower bound holds for exact determinantal complexity.
There is also a consequence for algebraic branching programs. An
affine-linear algebraic branching program is a finite directed acyclic
graph with a source and a sink and an affine-linear form on each edge; it
computes the sum of products of edge labels along all source-to-sink paths.
Over $\C$, computing $\per_m$ requires $\Omega(m^3)$ vertices and
$\Omega(m^3)$ edges. The same bounds hold for coefficientwise limits when
the relevant vertex or edge budget remains fixed along the approximation,
even if the graph varies. Section~\ref{sec:consequences} gives the precise
statements and determinant reductions.

\subsection{Earlier bounds and geometric methods}

The early matrix-order lower bounds were linear in the permanent order $m$.
For the variable-or-constant projection model, von zur Gathen published a
bound due to Babai and Seress with leading term $\sqrt{2}\,m$ and a loss of
$6\sqrt m$ \cite[Theorem~4.4]{vonZurGathen1987}. For arbitrary affine
representations over $\C$, Meshulam's final theorem
\cite[pp.~270--271]{Meshulam1989} and Cai's dimension argument
\cite[Theorem~1.5 and Section~2.2]{Cai1990} give the scale
$\sqrt{2}\,m-O(1)$. Their arguments compare the dimension of the affine
matrix image with the minimum rank that an exact permanent representation
must have.

Mignon and Ressayre raised the exact lower bound to
$\dc(\per_m)\ge m^2/2$ in characteristic zero
\cite{MignonRessayre2004}. Their proof compares Hessians at a zero:
affine pullback cannot increase Hessian rank, the Hessian of an order-$n$
determinant has rank at most $2n$ at a noninvertible matrix, and, for
$m\ge2$, the permanent has a zero where its Hessian has full rank $m^2$.
Cai, Chen and Li extended the quadratic scale to positive characteristic
different from two by constructing sparse zeros that avoid the large
factorial factors in the earlier Hessian calculation
\cite[Theorem~2.3]{CaiChenLi2010}.

On the upper side, Grenet's construction realizes $\per_m$ as a determinant
of order $2^m-1$, using a branching program indexed by the subsets of
$\{1,\ldots,m\}$ \cite[unnumbered Theorem, p.~3]{Grenet2012}.
Its determinant entries are individual variables and constants, so it is
also an upper bound in the unrestricted exact model.

Mulmuley and Sohoni's geometric complexity theory program studies
permanent--determinant separation through orbit closures of determinants
and padded permanents, seeking equations or representation-theoretic
obstructions to containment
\cite{MulmuleySohoni2001,MulmuleySohoni2008}.
Landsberg, Manivel and Ressayre proved the unrestricted complex border
bound $\bdc(\per_m)\ge m^2/2$ by studying hypersurfaces with degenerate
duals and the effect of padding
\cite[Theorem~1.1.1]{LandsbergManivelRessayre2010}.
Theorem~\ref{thm:main} raises the lower-bound exponent from two to three
in that same model. It remains a fixed polynomial bound. Whether the
permanent has superpolynomial determinantal complexity, or superpolynomial
border determinantal complexity, remains open.

Landsberg and Ressayre later proved exponential optima for exact
representations with symmetry requirements. Each specified projective
symmetry of the input must lift to a projective symmetry of the determinant
that fixes the pencil's constant matrix, preserves the image of its linear
part, and induces the given input action there. Over $\C$, for $m\ge3$,
requiring this for left multiplication by invertible diagonal and
permutation matrices gives minimum order $2^m-1$. Lifting the full
projective symmetry group of the permanent gives minimum order
$\binom{2m}{m}-1$ in the same range
\cite[Theorems~2.1 and~2.8; Propositions~2.9 and~2.10]{LandsbergRessayre2017}.
These hypotheses specify exact representation classes; the border model
of Theorem~\ref{thm:main} allows arbitrary affine determinant
approximations of fixed order.

Singularities also give geometric lower bounds for exact representations.
For a homogeneous polynomial of degree greater than two whose affine
singular locus has codimension greater than four, Alper, Bogart and Velasco
prove
$\dc(f)\ge\operatorname{codim}\operatorname{Sing}(f)+1$
\cite[Theorem~1.2]{AlperBogartVelasco2017}.
For the permanent, matrices with two zero columns show that this
codimension is at most $2m$, so this obstruction alone has only a linear
scale \cite[Remark~1.5]{AlperBogartVelasco2017}.
Our construction instead obtains a coefficient polynomial from a permanent
and finds a linear subspace on which its zero hypersurface is smooth. Its
degree is linear and the subspace dimension is quadratic in an auxiliary
matrix size.

For a smooth complex form of degree $r\ge2$ in $d\ge2$ variables, the classical
top polar degree is $r(r-1)^{d-2}$
\cite[Corollary~3.7]{Piene1978}.
The left and right kernel equations for the determinant matrix and the
Schur-complement reduction used below also appear in Sheshadri's June 2026
preprint \cite[Theorem~3 and Section~3.1]{SheshadriBorder2026}, as part of
its proposed conormal approach to border bounds for diagonal power sums.
Kumar and Volk's September 2026 preprint uses such kernel equations and a
multihomogeneous count to obtain a quadratic \emph{exact} lower bound for
$\sum_{i=1}^n x_i^n$
\cite[Theorem~1.1 and Section~3]{KumarVolk2026}.
Theorem~\ref{thm:smooth-border} proves the version needed here, including
persistence of finitely many simple polar intersections in one sufficiently
close determinant under an arbitrary coefficientwise approximation. Its
lower bound is proportional to the product $(r-1)(d-1)$.

\subsection{The coefficient construction and the proof}

To use this product at cubic scale, we seek an affine restriction of a
permanent of size proportional to $k$ that has degree proportional to $k$
in quadratically many variables and has a smooth projective zero set.
Let $k\ge1$ and $0\le q\le\lfloor k/2\rfloor$ be integers. With
indeterminates $\theta_1,\ldots,\theta_k$, put
$D=\diag(\theta_1,\ldots,\theta_k)$ and define
\[
 G_D(X)=[s^qt^q]\per_k(X+sI+tD)
 \in\mathbb Z[\theta_1,\ldots,\theta_k][X_{ij}].
\]
This polynomial is homogeneous of degree $r=k-2q$ in the $k^2$ entries
of $X$. We use the same notation after specializing the parameters in a
field. For the main construction we assume $k\ge128$ and take
$q=\lfloor(k-1)/3\rfloor$; then $r$ is linear in $k$ and
$T=\binom{q+2}{2}$ is quadratic in $k$.

The integral formula has two uses. Over $\C$, two count gates and a
cancellation gadget realize $G_D$, up to a fixed nonzero scalar, as an
affine restriction of a permanent of size $11k-2q$
(Proposition~\ref{prop:selector}). This holds for every complex choice of
the parameters. The scalar vanishes in characteristic two, so the geometric
argument specializes the integral coefficient formula itself. We will use
the complex projection at the complex parameter choice that geometry
eventually supplies.

Let $K=\overline{\mathbb F}_2$, choose pairwise distinct
$\bar\theta_1,\ldots,\bar\theta_k\in K$, and write
$\bar D=\diag(\bar\theta_1,\ldots,\bar\theta_k)$. Since permanent
and determinant agree over $K$, every matrix $X$ defines a plane curve
\[
 C_X=V\bigl(\det(sI+t\bar D+uX)\bigr)\subset\PP^2_K.
\]
Its fiber on $u=0$ consists of the distinct points
$[-\bar\theta_i:1:0]$. This fixed fiber makes the curve reduced and its
projection to $[t:u]$ finite and generically separable, and it controls the adjugate
of $sI+t\bar D+uX$ on the curve. The full adjugate defines a finite
subscheme $Z\subset C_X$. If all entrywise derivatives of $G_{\bar D}$
vanish at $X$, the coefficient $[s^qt^qu^{r-1}]$ of every adjugate entry
vanishes. We prove that these
entries span all degree-$(k-1)$ forms vanishing on $Z$, so the same
coefficient vanishes on every such form. For a degree-$q$ monomial
$s^it^ju^\ell$, multiplication by $s^{q-i}t^{q-j}u^{r-1-\ell}$ lets this
functional recover its coefficient in any degree-$q$ form. These exponents
are nonnegative because $r-1\ge q$. Thus restriction of degree-$q$ forms to $Z$ is
injective, giving the length bound in Proposition~\ref{prop:adjugate-length}:
\[
 K[s,t,u]_q\hookrightarrow H^0(Z,\cO_Z(q)),
 \qquad \length Z\ge\dim_K K[s,t,u]_q=T.
\]

The length of $Z$ is bounded above by the \emph{normalization defect}
$\delta(C_X)$, the length of the quotient between the normalization's
structure sheaf and the curve's structure sheaf. The comparison follows
the evaluation-ideal argument of Fantechi, G\"ottsche and van Straten
over $\C$ \cite[Proposition~C.2 and Corollary~C.3]{FGvS1999};
Section~\ref{sec:conductor} proves the needed conductor inclusion and
colength identity for our generically separable pencils in arbitrary
characteristic.
Thus every critical matrix lies over a curve with $\delta(C_X)\ge T$.

Applying the arbitrary-characteristic Severi dimension theorem of Christ,
He and Tyomkin \cite[Lemma~2.6, Proposition~2.7, and Remark~2.8]{CHT2023}
to the integral components bounds the locus of curves with this defect.
Their result builds on Tyomkin's dimension bound for toric pairs
\cite[Theorem~1.2(1)]{Tyomkin2013}. The sum
$R_X=(1,\ldots,1)\adj(sI+t\bar D+uX)$ of the adjugate rows controls
the matrices over one fixed curve: its zero subscheme has length
$k(k-1)/2$, and the common
boundary frame makes that subscheme determine the matrix uniquely.
A Hilbert-scheme estimate and an algebraic incidence then bound the
dimension of each matrix fiber. Combining the curve and matrix bounds gives
\[
 \dim\{X:\nabla G_{\bar D}(X)=0\}\le k^2+k-T,
 \qquad
 \codim\{X:\nabla G_{\bar D}(X)=0\}\ge T-k.
\]

The projective zero scheme of the raw partial derivatives forms a proper
family over the integral parameter space. Specialization transfers the
dimension bound to a nonempty open set of complex parameters. Over $\C$,
Euler's identity identifies these gradient zeros with the hypersurface's
singular scheme, and a general linear section avoids that scheme and meets
the smooth locus transversely. Proposition~\ref{prop:smooth-projection}
then supplies, for $k\ge128$ and the stated choice of $q$, a nonzero form
$G$ of degree $r$ in $d=\lfloor k^2/100\rfloor$ variables whose
projective zero set is smooth. For each $m\ge11k-2q$, it also supplies
a fixed affine map $\Phi_m$ satisfying $G=\per_m\circ\Phi_m$.

Figure~\ref{fig:proof-flow} shows how the two uses of the coefficient
formula meet and how the polar bound converts their output into the cubic
estimate.

\begin{figure}[tbp]
\centering
\setlength{\fboxsep}{7pt}
\fbox{\begin{minipage}{0.90\linewidth}
\centering
\textbf{One integral coefficient polynomial}\\[3pt]
$\displaystyle G_D(X)=[s^qt^q]\per_k(X+sI+tD)
\in\mathbb Z[\theta_1,\ldots,\theta_k][X_{ij}]$\\[3pt]
$\displaystyle D=\diag(\theta_1,\ldots,\theta_k),\qquad
r=k-2q,\qquad T=\binom{q+2}{2}$
\end{minipage}}

\smallskip
\begin{tabular}{@{}c@{\hspace{0.025\linewidth}}c@{}}
\begin{minipage}{0.43\linewidth}\centering
$\swarrow$\quad over $\C$, every parameter choice
\end{minipage}
&
\begin{minipage}{0.43\linewidth}\centering
$\searrow$\quad over $K=\overline{\mathbb F}_2$, choose distinct parameters
\end{minipage}
\\[4pt]
\fbox{\begin{minipage}{0.405\linewidth}\centering
\textbf{Permanent projection}\\[3pt]
$\displaystyle \per_m(A_m(X))=a_0G_D(X)$\\[3pt]
$m\ge11k-2q,\qquad a_0\ne0$
\end{minipage}}
&
\fbox{\begin{minipage}{0.405\linewidth}\centering
\textbf{Plane-curve geometry}\\[3pt]
$\displaystyle \nabla G_{\bar D}(X)=0
\ \Longrightarrow\ \delta(C_X)\ge T$\\[3pt]
$\displaystyle \codim V(\nabla G_{\bar D})\ge T-k$
\end{minipage}}
\\[5pt]
\begin{minipage}{0.43\linewidth}\centering
$\searrow$\quad apply at the complex parameters supplied by geometry
\end{minipage}
&
\begin{minipage}{0.43\linewidth}\centering
$\swarrow$\quad proper specialization and a transverse linear section
\end{minipage}
\end{tabular}

\smallskip
\fbox{\begin{minipage}{0.90\linewidth}
\centering
\textbf{A smooth affine restriction over $\C$}\\[3pt]
A nonzero homogeneous form $G$ of degree $r$ in $d$ variables has a smooth
projective zero set and satisfies $G=\per_m\circ\Phi_m$ for a fixed affine map $\Phi_m$.\\[3pt]
$\displaystyle r-1\ge\frac{k}{4},\qquad
d-1\ge\frac{k^2}{200}$
\end{minipage}}

\smallskip
$\big\downarrow$\quad the polar count for smooth homogeneous forms
\smallskip

\fbox{\begin{minipage}{0.90\linewidth}
\centering
$\displaystyle
\bdc(\per_m)\ge\bdc(G)
\ge\frac{(r-1)(d-1)}{4e}
\ge\frac{k^3}{3200e}$
\end{minipage}}
\caption{The two uses of the integral coefficient formula meet over
$\C$. The graph projection applies to every complex parameter choice;
the characteristic-two branch supplies a suitable complex choice through
the proper parameter family. The displayed scale bounds use
$q=\lfloor(k-1)/3\rfloor$, $k\ge128$, and $m\ge11k-2q$.}
\label{fig:proof-flow}
\end{figure}


Sections~\ref{sec:selector}--\ref{sec:transfer} establish this
construction. Section~\ref{sec:polar} then proves the smooth-form bound
stated in Section~\ref{sec:models}, using selected simple polar
intersections and the left and right kernels of one nearby determinant.
Section~\ref{sec:consequences} combines the two results and gives the
exact, branching-program, and smooth-initial-form consequences.

\FloatBarrier
